% !TEX root = ../projekt.tex
% =============================================================================
% STAV: hotové podľa aktuálneho kódu.
% =============================================================================
\chapter{Návod na spustenie simulátora}
\label{app:manual}

\section*{Preklad a spustenie}
\begin{lstlisting}
make                      # build rss_sim
./run.sh                  # simulation + table + plot, uses config.conf
./run.sh my_test.conf     # the same with another config
./rss_sim config.conf     # simulation only
\end{lstlisting}

\section*{Výstup}
Výsledkom je súbor CSV podľa kľúča \texttt{out} a graf s~rovnakým menom a príponou \texttt{.png}.
\begin{lstlisting}
method,window,num_channels,overload,fairness,chi
hash,1,8,110.914,0.922365,0.0851074
\end{lstlisting}
Riadky \texttt{hash} sú priemerom cez všetky kľúče. Orákulum \texttt{jsq} kľúče nepoužíva.

\TODO{Doplniť závislosti (kompilátor C++17, Python, pandas, matplotlib) a popis všetkých kľúčov.}
